% =============================================================================
\chapter{Markov Decision Processes}
\label{ch:mdp}
% =============================================================================

Imagine a robot learning to walk. At each moment, it observes its posture
(the \emph{state}), chooses a force to apply to its joints (the
\emph{action}), and receives a signal --- did it stay upright or fall over?
(the \emph{reward}). Its goal: find a strategy of actions that maximizes
rewards accumulated over time. This scheme --- state, action, reward, new
state --- is universal: it describes a chess player, an autonomous vehicle,
and a trading algorithm alike. The mathematical formalization of this scheme
has a name: the \emph{Markov decision process} (MDP), formalized by Richard
Bellman in the 1950s.

\begin{intuition}
A Markov Decision Process (MDP) is the central mathematical model of
reinforcement learning. It formalizes how an agent interacts with an
environment by making sequential decisions to maximize long-term cumulative
reward.
\end{intuition}

% -----------------------------------------------------------------------------
\section{Introduction}
% -----------------------------------------------------------------------------

Reinforcement learning rests on a simple idea: an agent learns to act by
interacting with an environment. At each time step $t$, the agent observes a
state $s_t \in \mathcal{S}$, chooses an action $a_t \in \mathcal{A}$, receives
a reward $r_{t+1} \in \R$, and transitions to a new state $s_{t+1}$. The
mathematical formalism that captures this interaction is the
\emph{Markov Decision Process}.

% -----------------------------------------------------------------------------
\section{Formal Definition}
% -----------------------------------------------------------------------------

\begin{definition}[Markov Decision Process]
An MDP is a five-tuple $\mathcal{M} = (\mathcal{S}, \mathcal{A}, P, R, \gamma)$ where:
\begin{itemize}
    \item $\mathcal{S}$ is a finite or countable set of \textbf{states},
    \item $\mathcal{A}$ is a finite or countable set of \textbf{actions},
    \item $P : \mathcal{S} \times \mathcal{A} \times \mathcal{S} \to [0,1]$ is the
    \textbf{transition function}: $P(s' \mid s, a) = \PP(S_{t+1} = s' \mid S_t = s, A_t = a)$,
    \item $R : \mathcal{S} \times \mathcal{A} \times \mathcal{S} \to \R$ is the
    \textbf{reward function}: $R(s, a, s')$ is the reward received when transitioning
    from $s$ to $s'$ under action $a$,
    \item $\gamma \in [0, 1)$ is the \textbf{discount factor}.
\end{itemize}
\end{definition}

\begin{keyformulas}[title=\textbf{Markov Property}]
The Markov property states that the future state depends only on the present
state and current action:
\[
    \PP(S_{t+1} = s' \mid S_0, A_0, S_1, A_1, \ldots, S_t, A_t)
    = \PP(S_{t+1} = s' \mid S_t, A_t) = P(s' \mid S_t, A_t).
\]
\end{keyformulas}

% -----------------------------------------------------------------------------
\section{Policy}
% -----------------------------------------------------------------------------

\begin{definition}[Policy]
A \textbf{policy} $\pi$ is a decision rule that prescribes the agent's behavior.
\begin{itemize}
    \item \textbf{Deterministic policy}: $\pi : \mathcal{S} \to \mathcal{A}$,
    where $a = \pi(s)$.
    \item \textbf{Stochastic policy}: $\pi : \mathcal{S} \times \mathcal{A} \to [0,1]$,
    where $\pi(a \mid s) = \PP(A_t = a \mid S_t = s)$ with
    $\sum_{a \in \mathcal{A}} \pi(a \mid s) = 1$ for all $s$.
\end{itemize}
\end{definition}

\begin{definition}[Stationary Policy]
A policy $\pi$ is \textbf{stationary} if it does not depend on time:
$\pi_t = \pi$ for all $t \geq 0$. Unless stated otherwise, all policies
considered in this course are stationary.
\end{definition}

% -----------------------------------------------------------------------------
\section{Value Functions}
% -----------------------------------------------------------------------------

\begin{definition}[Discounted Return]
The \textbf{discounted return} from time $t$ is:
\[
    G_t = \sum_{k=0}^{\infty} \gamma^k R_{t+k+1}.
\]
The factor $\gamma < 1$ ensures convergence of this sum when rewards are bounded.
\end{definition}

\begin{definition}[State-Value Function]
The \textbf{state-value function} under policy $\pi$ is:
\[
    V^\pi(s) = \E_\pi \left[ G_t \mid S_t = s \right]
    = \E_\pi \left[ \sum_{k=0}^{\infty} \gamma^k R_{t+k+1} \;\middle|\; S_t = s \right].
\]
\end{definition}

\begin{definition}[Action-Value Function]
The \textbf{action-value function} (or Q-function) under policy $\pi$ is:
\[
    Q^\pi(s, a) = \E_\pi \left[ G_t \mid S_t = s, A_t = a \right].
\]
\end{definition}

\begin{remark}
The relationship between $V^\pi$ and $Q^\pi$ is:
\[
    V^\pi(s) = \sum_{a \in \mathcal{A}} \pi(a \mid s) \, Q^\pi(s, a).
\]
\end{remark}

% -----------------------------------------------------------------------------
\section{Bellman Equations}
% -----------------------------------------------------------------------------

\begin{theorem}[Bellman Equation for $V^\pi$]
For any policy $\pi$ and any state $s \in \mathcal{S}$:
\[
    V^\pi(s) = \sum_{a \in \mathcal{A}} \pi(a \mid s) \sum_{s' \in \mathcal{S}}
    P(s' \mid s, a) \left[ R(s, a, s') + \gamma \, V^\pi(s') \right].
\]
\end{theorem}

\begin{proof}
By the definition of discounted return and the Markov property:
\begin{align*}
    V^\pi(s) &= \E_\pi[G_t \mid S_t = s] \\
    &= \E_\pi[R_{t+1} + \gamma G_{t+1} \mid S_t = s] \\
    &= \sum_{a} \pi(a \mid s) \sum_{s'} P(s' \mid s, a)
       \left[ R(s, a, s') + \gamma \, \E_\pi[G_{t+1} \mid S_{t+1} = s'] \right] \\
    &= \sum_{a} \pi(a \mid s) \sum_{s'} P(s' \mid s, a)
       \left[ R(s, a, s') + \gamma \, V^\pi(s') \right]. \qedhere
\end{align*}
\end{proof}

\begin{theorem}[Bellman Equation for $Q^\pi$]
For any policy $\pi$, any state $s$, and any action $a$:
\[
    Q^\pi(s, a) = \sum_{s' \in \mathcal{S}} P(s' \mid s, a)
    \left[ R(s, a, s') + \gamma \sum_{a' \in \mathcal{A}} \pi(a' \mid s') \, Q^\pi(s', a') \right].
\]
\end{theorem}

\begin{proof}
Expanding similarly:
\begin{align*}
    Q^\pi(s, a) &= \E_\pi[R_{t+1} + \gamma G_{t+1} \mid S_t = s, A_t = a] \\
    &= \sum_{s'} P(s' \mid s, a) \left[ R(s, a, s') +
       \gamma \, \E_\pi[G_{t+1} \mid S_{t+1} = s'] \right] \\
    &= \sum_{s'} P(s' \mid s, a) \left[ R(s, a, s') +
       \gamma \sum_{a'} \pi(a' \mid s') Q^\pi(s', a') \right]. \qedhere
\end{align*}
\end{proof}

\begin{keyformulas}[title=\textbf{Bellman Equations --- Summary}]
\begin{align}
    V^\pi(s) &= \sum_a \pi(a|s) \sum_{s'} P(s'|s,a)[R(s,a,s') + \gamma V^\pi(s')]
    \label{eq:bellman_v} \\
    Q^\pi(s,a) &= \sum_{s'} P(s'|s,a)[R(s,a,s') + \gamma \sum_{a'} \pi(a'|s') Q^\pi(s',a')]
    \label{eq:bellman_q}
\end{align}
\end{keyformulas}

% -----------------------------------------------------------------------------
\section{Optimal Policy and Bellman Optimality Equations}
% -----------------------------------------------------------------------------

\begin{definition}[Optimal Policy]
A policy $\pi^*$ is \textbf{optimal} if for every state $s \in \mathcal{S}$:
\[
    V^{\pi^*}(s) \geq V^\pi(s) \quad \text{for all policies } \pi.
\]
We write $V^*(s) = V^{\pi^*}(s)$ and $Q^*(s,a) = Q^{\pi^*}(s,a)$.
\end{definition}

\begin{theorem}[Bellman Optimality Equations]
The optimal value functions satisfy:
\begin{align}
    V^*(s) &= \max_{a \in \mathcal{A}} \sum_{s'} P(s'|s,a)
    \left[ R(s,a,s') + \gamma V^*(s') \right] \label{eq:bellman_opt_v} \\
    Q^*(s,a) &= \sum_{s'} P(s'|s,a)
    \left[ R(s,a,s') + \gamma \max_{a'} Q^*(s',a') \right] \label{eq:bellman_opt_q}
\end{align}
\end{theorem}

\begin{proof}
For Equation~\eqref{eq:bellman_opt_v}, we use the fact that the optimal policy
selects the action maximizing value:
\begin{align*}
    V^*(s) &= \max_\pi V^\pi(s) = \max_a Q^*(s,a)
    = \max_a \sum_{s'} P(s'|s,a)[R(s,a,s') + \gamma V^*(s')].
\end{align*}
For Equation~\eqref{eq:bellman_opt_q}, we substitute $V^*(s') = \max_{a'} Q^*(s',a')$
into the Bellman equation for $Q^*$.
\end{proof}

\begin{theorem}[Existence of a Deterministic Optimal Policy]
\label{thm:existence_opt}
For any finite MDP, there exists at least one deterministic optimal policy
$\pi^* : \mathcal{S} \to \mathcal{A}$ defined by:
\[
    \pi^*(s) = \argmax_{a \in \mathcal{A}} Q^*(s, a).
\]
\end{theorem}

% -----------------------------------------------------------------------------
\section{Bellman Operator and Contraction}
% -----------------------------------------------------------------------------

\begin{definition}[Bellman Operator]
The \textbf{Bellman operator} $\mathcal{T}^\pi$ for a policy $\pi$ is defined by:
\[
    (\mathcal{T}^\pi V)(s) = \sum_a \pi(a|s) \sum_{s'} P(s'|s,a)
    [R(s,a,s') + \gamma V(s')].
\]
The \textbf{Bellman optimality operator} is:
\[
    (\mathcal{T}^* V)(s) = \max_a \sum_{s'} P(s'|s,a)
    [R(s,a,s') + \gamma V(s')].
\]
\end{definition}

\begin{theorem}[Contraction of the Bellman Operator]
\label{thm:contraction}
The operators $\mathcal{T}^\pi$ and $\mathcal{T}^*$ are contractions in the
infinity norm with factor $\gamma$:
\[
    \norm{\mathcal{T}^\pi V - \mathcal{T}^\pi U}_\infty
    \leq \gamma \norm{V - U}_\infty.
\]
By Banach's fixed-point theorem, each has a unique fixed point:
$V^\pi$ for $\mathcal{T}^\pi$ and $V^*$ for $\mathcal{T}^*$.
\end{theorem}

\begin{proof}
For any $s \in \mathcal{S}$:
\begin{align*}
    \abs{(\mathcal{T}^* V)(s) - (\mathcal{T}^* U)(s)}
    &= \abs{\max_a \sum_{s'} P(s'|s,a)[R + \gamma V(s')]
       - \max_a \sum_{s'} P(s'|s,a)[R + \gamma U(s')]} \\
    &\leq \max_a \sum_{s'} P(s'|s,a) \gamma \abs{V(s') - U(s')} \\
    &\leq \gamma \norm{V - U}_\infty.
\end{align*}
The transition from the second to third line uses
$\abs{\max_a f(a) - \max_a g(a)} \leq \max_a \abs{f(a) - g(a)}$.
\end{proof}

% -----------------------------------------------------------------------------
\section{Advantage Function}
% -----------------------------------------------------------------------------

\begin{definition}[Advantage Function]
The \textbf{advantage function} under policy $\pi$ is:
\[
    A^\pi(s, a) = Q^\pi(s, a) - V^\pi(s).
\]
It measures the relative advantage of taking action $a$ in state $s$ compared
to the average behavior under $\pi$.
\end{definition}

\begin{proposition}
For any policy $\pi$ and any state $s$:
\[
    \sum_{a \in \mathcal{A}} \pi(a \mid s) \, A^\pi(s, a) = 0.
\]
\end{proposition}

\begin{proof}
$\sum_a \pi(a|s) A^\pi(s,a) = \sum_a \pi(a|s)[Q^\pi(s,a) - V^\pi(s)]
= V^\pi(s) - V^\pi(s) = 0$.
\end{proof}

% -----------------------------------------------------------------------------
\section{MDP Examples}
% -----------------------------------------------------------------------------

\begin{example}[Gridworld]
Consider a $4 \times 4$ grid where an agent moves in four cardinal directions.
The terminal state is cell $(4,4)$ with reward $+1$. Each step costs $-0.01$.
Moves against walls leave the agent in place. This formalizes as an MDP with:
\begin{itemize}
    \item $\mathcal{S} = \{(i,j) : 1 \leq i,j \leq 4\}$, $|\mathcal{S}| = 16$,
    \item $\mathcal{A} = \{\text{up, down, left, right}\}$,
    \item Deterministic transitions (except at boundaries),
    \item $\gamma = 0.99$.
\end{itemize}
\end{example}

\begin{example}[$K$-Armed Bandit]
A $K$-armed bandit is a degenerate MDP with $|\mathcal{S}| = 1$ (single state).
The agent selects an arm $a \in \{1, \ldots, K\}$ at each step and receives a
random reward $R_a \sim \mathcal{D}_a$. This special case allows studying the
exploration-exploitation dilemma in isolation.
\end{example}

% -----------------------------------------------------------------------------
\section{Python Implementation}
% -----------------------------------------------------------------------------

\begin{codeblock}[title=\textbf{Defining a Simple MDP with Gymnasium}]
\begin{minted}{python}
import numpy as np
import gymnasium as gym
from gymnasium import spaces

class GridWorldEnv(gym.Env):
    """Custom 4x4 Gridworld environment."""
    metadata = {"render_modes": ["human"]}

    def __init__(self, size=4, gamma=0.99):
        super().__init__()
        self.size = size
        self.gamma = gamma
        self.observation_space = spaces.Discrete(size * size)
        self.action_space = spaces.Discrete(4)  # up, down, left, right
        self.goal = (size - 1, size - 1)
        self._action_to_dir = {
            0: np.array([-1, 0]),  # up
            1: np.array([1, 0]),   # down
            2: np.array([0, -1]), # left
            3: np.array([0, 1]),  # right
        }
        self.reset()

    def _pos_to_state(self, pos):
        return pos[0] * self.size + pos[1]

    def reset(self, seed=None, options=None):
        super().reset(seed=seed)
        self.agent_pos = np.array([0, 0])
        return self._pos_to_state(self.agent_pos), {}

    def step(self, action):
        direction = self._action_to_dir[action]
        new_pos = np.clip(self.agent_pos + direction, 0, self.size - 1)
        self.agent_pos = new_pos
        terminated = tuple(self.agent_pos) == self.goal
        reward = 1.0 if terminated else -0.01
        return self._pos_to_state(self.agent_pos), reward, terminated, False, {}
\end{minted}
\end{codeblock}

\begin{codeblock}[title=\textbf{Building the Transition Matrix}]
\begin{minted}{python}
def build_transition_matrix(env):
    """Build P[s, a, s'] and R[s, a, s'] for a finite MDP."""
    nS = env.observation_space.n
    nA = env.action_space.n
    P = np.zeros((nS, nA, nS))
    R = np.zeros((nS, nA, nS))

    for s in range(nS):
        for a in range(nA):
            env.agent_pos = np.array([s // env.size, s % env.size])
            s_next, reward, _, _, _ = env.step(a)
            P[s, a, s_next] = 1.0
            R[s, a, s_next] = reward
    return P, R

env = GridWorldEnv(size=4)
P, R = build_transition_matrix(env)
print(f"P shape: {P.shape}")  # (16, 4, 16)
print(f"R shape: {R.shape}")  # (16, 4, 16)
\end{minted}
\end{codeblock}

% -----------------------------------------------------------------------------
\section{Direct Algebraic Solution}
% -----------------------------------------------------------------------------

\begin{proposition}[Matrix Solution of the Bellman Equation]
For a finite MDP with $|\mathcal{S}| = n$, the Bellman equation for $V^\pi$
can be written in matrix form:
\[
    \mathbf{v}^\pi = \mathbf{r}^\pi + \gamma \, \mathbf{P}^\pi \mathbf{v}^\pi
\]
where $\mathbf{P}^\pi_{s,s'} = \sum_a \pi(a|s) P(s'|s,a)$ and
$\mathbf{r}^\pi_s = \sum_a \pi(a|s) \sum_{s'} P(s'|s,a) R(s,a,s')$.
The solution is:
\[
    \mathbf{v}^\pi = (I - \gamma \mathbf{P}^\pi)^{-1} \mathbf{r}^\pi.
\]
\end{proposition}

\begin{warning}[title=\textbf{Complexity of Direct Solution}]
Matrix inversion has complexity $O(n^3)$, making it impractical for MDPs with
large state spaces. This motivates the iterative methods of
Chapter~\ref{ch:dynamic_programming}.
\end{warning}

% -----------------------------------------------------------------------------
\section{Extensions of the MDP Model}
% -----------------------------------------------------------------------------

\begin{definition}[POMDP]
A \textbf{Partially Observable MDP} (POMDP) augments an MDP with an observation
set $\mathcal{O}$ and an observation function $O : \mathcal{S} \times \mathcal{A} \to \Delta(\mathcal{O})$.
The agent does not directly observe state $s_t$ but instead receives observation
$o_t \sim O(\cdot | s_t, a_{t-1})$.
\end{definition}

\begin{definition}[Finite-Horizon MDP]
A \textbf{finite-horizon} MDP with horizon $H$ terminates after exactly $H$
time steps. The return is:
\[
    G_t = \sum_{k=0}^{H-t-1} \gamma^k R_{t+k+1}.
\]
\end{definition}

\begin{definition}[Continuous MDP]
When $\mathcal{S} \subseteq \R^n$ and/or $\mathcal{A} \subseteq \R^m$, we speak
of continuous state or action space MDPs. Sums are replaced by integrals in the
Bellman equations.
\end{definition}

% -----------------------------------------------------------------------------
\section{Exercises}
% -----------------------------------------------------------------------------

\begin{exercise}[Bellman Equation]
Consider an MDP with 3 states $\{s_1, s_2, s_3\}$, $\gamma = 0.9$, and the
uniform policy $\pi(a|s) = 1/2$ for two actions. The transitions and rewards are:
\begin{itemize}
    \item From $s_1$, action $a_1$: to $s_2$ with $R=1$; action $a_2$: to $s_3$ with $R=0$.
    \item From $s_2$, any action: to $s_1$ with $R=2$.
    \item From $s_3$, any action: to $s_3$ with $R=0$ (absorbing state).
\end{itemize}
Write the system of Bellman equations for $V^\pi$ and solve it.
\end{exercise}

\begin{exercise}[Contraction Operator]
Show that if $\gamma = 0$, the Bellman optimality operator converges in a single
iteration. What is the interpretation of $\gamma = 0$?
\end{exercise}

\begin{exercise}[Implementation]
Implement a custom \texttt{FrozenLake} environment with Gymnasium. Build the
transition matrix and compute $V^\pi$ by matrix inversion for the uniform policy.
Verify your result by comparing with Gymnasium's \texttt{FrozenLake-v1}.
\end{exercise}

\begin{exercise}[Advantage Function]
Show that for a deterministic optimal policy $\pi^*$:
\[
    A^{\pi^*}(s, \pi^*(s)) = 0 \quad \text{and} \quad A^{\pi^*}(s, a) \leq 0
    \quad \forall a \neq \pi^*(s).
\]
\end{exercise}
